\chapter*{Заключение}                       % Заголовок
\addcontentsline{toc}{chapter}{Заключение}  % Добавляем его в оглавление

В~ходе выполненного исследования решена научно-техническая задача,
связанная с~разработкой и~верификацией методики распознавания
и~реидентификации человека, интегрирующей детектирование силуэтов
и~лиц с~модульной системой признаковых представлений и~их вероятностной
интеграцией. На~основе обзора существующих методов и~алгоритмов
детекции объектов, реидентификации и~комбинирования эмбеддингов
сформулирован ряд ограничений современных систем, в~частности
недостаточная точность идентификации в~сложных условиях съемки,
зависимость от~одного типа признаков и~неэффективность при высокой
вариативности внешнего облика человека. Актуальность темы
и~поставленная цель исследования, заключающаяся в~разработке методики,
обеспечивающей повышение точности распознавания и~реидентификации
людей при рациональных вычислительных затратах, обусловили
необходимость комплексной методики, объединяющей в~едином контуре
согласованное детектирование связанных компонентов <<силуэт-лицо>>,
формирование их дискриминативных эмбеддингов и~вероятностную интеграцию
модальностей при принятии решения.

\textbf{Основные результаты исследования заключаются в~следующем.}
%% Согласно ГОСТ Р 7.0.11-2011:
%% 5.3.3 В заключении диссертации излагают итоги выполненного исследования, рекомендации, перспективы дальнейшей разработки темы.
%% 9.2.3 В заключении автореферата диссертации излагают итоги данного исследования, рекомендации и перспективы дальнейшей разработки темы.
\begin{enumerate}
  \item Разработана комплексная методика распознавания многокомпонентных
        структур по~их разрозненным частям, реализующая единый конвейер
        обработки из~четырех последовательных этапов: (1)~согласованное
        детектирование связанных компонентов структуры; (2)~раздельное
        извлечение дискриминативных эмбеддингов компонентов кодировщиками
        Vision Transformer с~функцией потерь ArcFace; (3)~вычисление
        гауссовых правдоподобий для каждого класса; (4)~принятие решения
        по~многоклассовому MAP-критерию. Методика обеспечивает поэтапное
        повышение качества на~каждой стадии при минимизации каскадного
        эффекта накопления ошибок и~устойчивое распознавание при неполноте
        данных; апробирована на~наиболее представительном частном случае~--
        реидентификации людей в~информационно-поисковых системах по~паре
        связанных компонентов~<<силуэт-лицо>> ($K=2$).

  \item Разработана байесовская модель объединения признаковых описаний
        компонентов многокомпонентной структуры и~решающее правило
        по~многоклассовому MAP-критерию с~автоматическим взвешиванием
        компонентов через ковариационные матрицы класса $\Sigma_y$
        в~гауссовых правдоподобиях. Модель обеспечивает адаптивный вклад
        каждого компонента без ручной подстройки весов и~корректное
        поведение решающего правила при произвольном множестве доступных
        компонентов (временном отсутствии или зашумлении части из~них).
        На~наиболее представительном частном случае реидентификации
        личности (набор CUHK03) достигнуты макроусредненные значения:
        Precision~-- 95,65\,\%, Recall~-- 95,54\,\%, $F_1$-мера~-- 94,31\,\%,
        Accuracy~-- 93,42\,\%; по~метрике mAP получены значения 95,65\,\% /
        98,3\,\% / 89,2\,\% для наборов CUHK03 / Market-1501 / MARS
        соответственно, что превосходит базовые способы линейного
        объединения признаков.

  \item Разработана модификация критерия обучения детектора (на~базе
        архитектуры YOLOv8) для согласованной локализации связанных
        компонентов структуры (внешнего и~вложенного), основанная
        на~введении дополнительного компонента функции потерь
        $\mathcal{L}_{\text{inner}}$ и~управляющего коэффициента
        $\lambda_{\text{inner}}$, совместно штрафующих ошибки локализации
        вложенного компонента в~границах внешнего. Экспериментально
        показано, что значения $\lambda_{\text{inner}}$ в~диапазоне
        2,0--2,7 обеспечивают наилучшие показатели по~детекции вложенного
        компонента (в~случае $K=2$~-- лица в~границах силуэта) без
        статистически значимой деградации качества детектирования базового
        внешнего класса.

  \item Подтверждена практическая значимость предложенной методики.
        Результаты работы внедрены и~используются в~деятельности
        ФГБУ~<<НИЦ~<<Институт имени Н.\,Е.\,Жуковского>>>>
        (сокращение времени принятия решения до~5\,\% в~системах
        визуального наблюдения).
\end{enumerate}


Проведенные в~третьей главе эксперименты показали, что модель YOLOv8
с~измененной функцией потерь повышает точность детекции
лиц при сохранении стабильных показателей по~классу~<<человек>>,
особенно в~сложных условиях частичного закрытия лица или неустойчивого
освещения. Оптимальные значения $\lambda_{\text{inner}}$ в~диапазоне
2,0--2,7 обеспечили увеличение mAP@0.5 примерно на~0,9 процентного
пункта (с~0,714 до~0,723) по~сравнению с~базовой реализацией, не~приводя к~ухудшению результатов
для других объектов. Анализ ViT-эмбеддингов, сформированных для
набора данных CUHK03, продемонстрировал образование четких кластеров,
соответствующих конкретным личностям и~подтверждающих
высокую дискриминационную способность представленных векторов.
Количественные метрики, рассчитанные для эмбеддингов лиц и~силуэтов,
указывают на~низкий уровень ложных срабатываний и~отказов даже при
жестких порогах. Практическая значимость подтверждается
экспериментами, в~которых достигнута повышенная точность реидентификации
при рациональных вычислительных затратах.

\textbf{Научная новизна} результатов работы выражается в~разработке
оригинальной методики, сочетающей согласованное детектирование
силуэтов и~лиц с~раздельным построением дискриминативных эмбеддингов
для каждой модальности, а~также в~создании нового метода вероятностной
интеграции лицевых и~силуэтных признаков на~основе байесовской модели
и~многоклассового MAP-критерия с~автоматическим взвешиванием
модальностей через ковариационные матрицы класса. Впервые предложенная
комплексная методика обеспечивает интегрированный учет лицевых
и~силуэтных признаков с~корректным поведением решающего правила при
неполноте данных, и~обеспечивает повышение точности реидентификации
в~режиме, близком к~реальному времени, при высоком уровне
вариативности исходных данных.

\textbf{Перспективы развития} результатов исследования заключаются
в~расширении набора признаков, учитываемых в~формировании итогового
эмбеддинга (в~частности, атрибутов одежды, походки, характерных
аксессуаров), а~также в~применении моделей, использующих
дополнительные архитектурные новшества (гибридные решения на~основе
сверточных нейронных сетей и~трансформеров, усовершенствованные
механизмы самовнимания). Перспективным направлением является
исследование смесей гауссовых распределений и~непараметрических
методов оценки плотностей для повышения гибкости вероятностной модели.
Кроме того, целесообразно исследовать возможности распределенного
обучения и~более глубокого переноса (transfer learning), что
востребовано при развертывании систем реидентификации в~рамках
гетерогенных сред и~систем интеллектуального видеонаблюдения.

Результаты диссертации свидетельствуют о~том, что
разработанная методика распознавания и~реидентификации человека при
интеграции в~информационно-поисковые системы обеспечивает повышение
общей точности реидентификации и~сокращение числа ложных тревог при
сохранении надежности и~детальности анализа в~режиме реального
времени. Полученные выводы и~рекомендации имеют как теоретическое,
так и~прикладное значение, обеспечивая возможность
разработки и~адаптации систем распознавания для решения различных
задач в~области безопасности, мониторинга и~управления доступом.

В~заключение автор выражает благодарность и~признательность научному
руководителю Гончаренко~В.\,И. за~поддержку, помощь, обсуждение
результатов и~научное руководство, а~также сотрудникам Института
проблем управления им.\,В.\,А.\,Трапезникова Российской академии наук
за~ценные замечания и~рекомендации на~разных этапах подготовки
диссертации.
